\section{おわりに}
本稿では，物体領域の検出とカテゴリのアノテーションを分離し，WordNetの下位語をたどることで追加学習なしに粒度を変更できる物体アノテーションを，DETRとSigLIPによる提案手法，YOLOの予測を深さ0とするハイブリッド手法，およびYOLO単体で比較した．
COCO 2017 Validationでは，ハイブリッド手法が深さ0のTop-1正解率0.969をYOLO11xと同じく保ちつつ，粒度を変更できたインスタンスを提案手法の1.79倍に増やした．
LVISの細かい粒度の正解による評価では，シグモイドの累積確信度による停止規則を用いたハイブリッド手法の階層的F値が0.659となり，YOLO11x単体の0.604を上回った．
COCO外の物体については，COCOで学習したYOLO11xとDETRは閾値を下げても再現率が2割未満で飽和する一方，OWLv2の物体らしさのスコアによる候補をYOLO11xのボックスと重ならない領域に限って加えることで，COCOの物体の精度を保ったまま，COCO外の物体の検出を4.6\%から31.6\%に増やし，粒度の変更の対象にできることを示した．
一方，検出したCOCO外の物体のうち深さ0を正しく命名できたのは8.1\%であり，語彙の選択と不適切な領域の棄却が今後の課題である．
\add{これに対し，物体らしさとシグモイドの確信度による候補の棄却，頻出するカテゴリの語から順に下位語を選ぶ命名，およびSigLIP SO400Mを組み合わせることで，1画像あたり5件の候補を加える条件で，COCO外の物体の命名を1.7\%から4.1\%に，加えた候補の命名の適合率を3.6\%から8.1\%に向上させた．}
